\scanpage{19}
\mspage{4)}
Récapitulation (II)

$k$ corps \uline{alg. clos}, car. 0 i.e. \deleted{\unsure{k}} $\mathbb Q\longrightarrow k$,
$X$ courbe de type $g,\nu$ sur $k$, $x\in X$ ; d'où
\[
\xi\in\deleted{\unsure{\mathrm{Pt}}}M_{g,\nu}(k)=\uline{\Gamma}(M_{g,\nu\,k}/k)
\qquad X=\mathcal X_{g,\nu\,\xi}=(\mathcal X_{g,\nu\,k})_\xi
\]
\[x\in\mathcal X_{g,\nu}(k)\]
$\xi,x$ pts géométriques de $M_{g,\nu\,\mathbb Q},\mathcal X_{g,\nu\,\mathbb Q}$,
d'où par fibration double
\[
(0)\quad\left\{\diagram{\begin{tikzcd}
\mathcal X_{g,\nu\,\mathbb Q}\arrow[r,no head]\arrow[d]&\mathcal X_{g,\nu\,k}\arrow[d,no head]\\
M_{g,\nu\,\mathbb Q}\arrow[d,no head]&M_{g,\nu\,k}\arrow[d,no head]\\
\mathbb Q\arrow[r,no head]&k
\end{tikzcd}}\right.
\]
le diagramme de suites exactes
\diagram{\begin{tikzcd}
&1\arrow[d]&1\arrow[d]&&\\
&\Pi_{X,x}\arrow[r,Rightarrow]\arrow[d]&\unsure{\Gamma_{X,x}|\Pi_{X,x}}\arrow[d]\arrow[r,no head]&{}&\\
1\arrow[r]&\mathcal G_{X,x}\arrow[r]\arrow[d]&\mathcal M_{X,x}\arrow[r]\arrow[d]&\Gamma_{\widetilde{\mathbb Q}/\mathbb Q}\arrow[r,no head]\arrow[d,equal]&1\\
1\arrow[r]&\Gamma_X\arrow[r]\arrow[d]&\mathcal N_{X\deleted{,x}}\arrow[r,no head]\arrow[d]&\Gamma_{\widetilde{\mathbb Q}/\mathbb Q}\arrow[r,no head]\arrow[d,no head]&1\\
&1&1&\unsure{1}&
\end{tikzcd}}
\margin{(1)}
\margin{$\widetilde{\mathbb Q}$ est la clôture alg. de $\mathbb Q$ dans $k$}
\margin{\uline{NB} Ce diagramme équivaut à la donnée de
\[
1\subset\overbrace{\underbrace{\Pi_{X,x}\subset\mathcal G_{X,x}}_{\Gamma_X}
\underbrace{{}\subset\mathcal M_{X,x}}_{\deleted{\unsure{\mathcal N}}\Gamma_{\widetilde{\mathbb Q}/\mathbb Q}}}^{\mathcal N_X}
\]}
où \deleted{$\unsure{\mathcal G_{X,x}\text{ est}}$}
\[
(2)\quad\mathrm{déf}\left\{\begin{aligned}
\Pi_{X,x}&=\pi_1(X,x) &&(\simeq\widehat\Pi_{g,\nu}\text{ non cant.})\\
\unsure{\mathcal G_{X,x}}&=\pi_1(\mathcal X_{g,\nu\,k},x)&&(\simeq\widehat{\mathcal G}_{g,\nu}\text{ ---})\\
\Gamma_X&=\pi_1(M_{g,\nu\,k},\xi)&&(\simeq\widehat\Gamma_{g,\nu}\ \ldots)\\
\mathcal M_{X,x}&=\pi_1(\unsure{\mathcal X_{g,\nu}|M_{g,\nu}},\deleted{\xi}\inserted{x})&&(\simeq\mathcal M_{g,\nu}\ \ldots)\\
\mathcal N_X&=\pi_1(M_{g,\nu},\xi)&&(\simeq\mathcal N_{g,\nu}\ \ldots)
\end{aligned}\right.
\]
Les deuxièmes mb $\widehat\Pi_{g,\nu}$ se définissent comme les premiers,
mais à l'aide de $(X_0,x_0)$ sur $\mathbb C$, de type $g,\nu$,
définissant les pts géom. $x_0$ et $\xi_0$ de
$\mathcal X_{g,\nu\,\mathbb Q}$ et $M_{g,\nu\,\mathbb Q}$.
\mspage{18}
